% Part: sets-functions-relations
% Chapter: functions
% Section: partial-functions

\documentclass[../../../include/open-logic-section]{subfiles}

\begin{document}

\olfileid{sfr}{fun}{par}

\olsection{పాక్షిక ప్రమేయాలు}

\begin{explain}
సాధ్యమైన అన్ని ఇన్‌పుట్లకూ ఫలితం నిర్వచించబడి ఉండాలనే నిబంధనను
సడలించడం కొన్నిసార్లు ఉపయోగకరం. ఇటువంటి ప్రతిచిత్రణాలను
\emph{పాక్షిక ప్రమేయాలు} అంటాం.
\end{explain}

\begin{defn}
$A$లోని ప్రతి \tetoken{మూలకానికి}{!!{element}}, $B$లో గరిష్ఠంగా ఒక
\tetoken{మూలకాన్ని}{!!{element}} కేటాయించే ప్రతిచిత్రణాన్ని
\emph{పాక్షిక ప్రమేయం} $f \colon A \pto B$ అంటాం.
$x \in A$కు $B$లోని ఒక మూలకాన్ని $f$ కేటాయిస్తే, $f(x)$
\emph{నిర్వచితం} అంటాం; లేకపోతే \emph{అనిర్వచితం} అంటాం.
$f(x)$ నిర్వచితమైతే $f(x) \fdefined$ అనీ, లేకపోతే
$f(x) \fundefined$ అనీ రాస్తాం. పాక్షిక ప్రమేయం $f$ యొక్క
\emph{ప్రవేశం}, అది నిర్వచితమయ్యే $A$ యొక్క ఉపసమితి:
$\dom{f} = \Setabs{x \in A}{f(x) \fdefined}$.
\end{defn}

\begin{ex}
ప్రతి ప్రమేయం $f\colon A \to B$ ఒక పాక్షిక ప్రమేయం కూడా అవుతుంది.
$A$ అంతటా నిర్వచితమయ్యే పాక్షిక ప్రమేయాలను---అంటే ఇప్పటివరకు మనం
కేవలం ప్రమేయాలుగా పిలిచినవాటిని---\emph{సర్వత్ర నిర్వచిత}
ప్రమేయాలు అని కూడా అంటాం.
\end{ex}

\begin{ex}
$f(x) = 1/x$ ద్వారా ఇచ్చే పాక్షిక ప్రమేయం $f \colon \Real \pto \Real$,
$x = 0$ వద్ద అనిర్వచితం; మిగిలిన అన్ని చోట్లా నిర్వచితం.
\end{ex}

\begin{prob}
$f\colon A \pto B$ ఇచ్చినప్పుడు, $g\colon B \pto A$ అనే పాక్షిక
ప్రమేయాన్ని ఇలా నిర్వచించండి: ఏ $y \in B$కైనా, $f(x) = y$ అయ్యే
ఏకైక $x \in A$ ఉంటే $g(y) = x$; లేకపోతే $g(y) \fundefined$.
$f$ ఒకటి-ఒకటి అయితే, ప్రతి $x \in \dom{f}$కూ $g(f(x)) = x$ అనీ,
ప్రతి $y \in \ran{f}$కూ $f(g(y)) = y$ అనీ చూపండి.
\end{prob}

\begin{defn}[పాక్షిక ప్రమేయపు గ్రాఫు]
$f\colon A \pto B$ ఒక పాక్షిక ప్రమేయం అనుకుందాం.
$f$ యొక్క \emph{గ్రాఫు} అంటే కింది విధంగా నిర్వచించే
$R_f \subseteq A \times B$ అనే సంబంధం:
\[
R_f = \Setabs{\tuple{x,y}}{f(x) = y}.
\]
\end{defn}

\begin{prop}
$R \subseteq A \times B$కు ఈ ధర్మం ఉందనుకుందాం:
$Rxy$, $Rxy'$ అయినప్పుడల్లా $y = y'$.
అప్పుడు $R$ అనేది కింది విధంగా నిర్వచించే పాక్షిక ప్రమేయం
$f\colon A \pto B$ యొక్క గ్రాఫు: $Rxy$ అయ్యే $y$ ఉంటే $f(x) = y$;
లేకపోతే $f(x) \fundefined$. $R$ అదనంగా \emph{సీరియల్} కూడా అయితే,
అంటే ప్రతి $x \in A$కూ $Rxy$ అయ్యే $y \in B$ ఉంటే,
$f$ సర్వత్ర నిర్వచితమవుతుంది.
\end{prop}

\begin{proof}
$Rxy$ అయ్యే $y$ ఉందనుకుందాం. $Rxy'$ అయ్యే మరో $y' \neq y$ ఉంటే,
$R$పై విధించిన షరతు ఉల్లంఘించబడుతుంది. కనుక $Rxy$ అయ్యే $y$ ఉంటే,
ఆ $y$ ఏకైకం; అందువల్ల $f$ సునిర్వచితం.
$R_f = R$ అనీ, $R$ సీరియల్ అయితే $f$ సర్వత్ర నిర్వచితం అనీ స్పష్టమే.
\end{proof}

\end{document}
